% BHU PYQ -- Professional Exam Template
\documentclass[11pt,a4paper]{article}

\usepackage[a4paper,top=2.5cm,bottom=2.5cm,left=2.8cm,right=2.8cm,headheight=14pt]{geometry}
\usepackage{amsmath,amssymb}
\usepackage[T1]{fontenc}
\usepackage[utf8]{inputenc}
\usepackage{lmodern}
\usepackage{microtype}
\usepackage{fancyhdr}
\usepackage{enumitem}
\usepackage{listings}
\usepackage{hyperref}
\usepackage{lastpage}
\usepackage{parskip}

\hypersetup{colorlinks=true,urlcolor=black,linkcolor=black,
  pdftitle={BVCA-201: Introduction To Databases -- BHU PYQ}}

\pagestyle{fancy}
\fancyhf{}
\renewcommand{\headrulewidth}{0.4pt}
\renewcommand{\footrulewidth}{0.4pt}
\fancyhead[L]{\small\textit{Banaras Hindu University}}
\fancyhead[R]{\small\textit{BVCA-201: Introduction To Databases}}
\fancyfoot[C]{\small Page \thepage\ of \pageref{LastPage}}

\lstset{
  basicstyle=\small\ttfamily,
  frame=single,
  breaklines=true,
  columns=flexible,
  keepspaces=true
}

\newlist{parts}{enumerate}{2}
\setlist[parts,1]{label=(\alph*),leftmargin=*,itemsep=4pt,topsep=3pt}
\setlist[parts,2]{label=(\roman*),leftmargin=*,itemsep=2pt,topsep=2pt}
\newcommand{\pts}[1]{\hfill\textit{[#1]}}

\begin{document}
\begin{center}
  {\large\bfseries BANARAS HINDU UNIVERSITY}\\[2pt]
  {\normalsize B.Voc. Semester II Examination, 2021-22}\\[6pt]
  \rule{0.8\linewidth}{0.5pt}\\[6pt]
  {\large\bfseries Paper: BVCA-201 --- Introduction To Databases}\\[4pt]
  {\normalsize Time: Three Hours \hfill Maximum Marks: 70}
\end{center}

\rule{\linewidth}{0.4pt}

\smallskip
\noindent\textit{Instructions: Answer the following (each question carry}

\bigskip

REFNO 052213
— Paper : BVCA-201
| Introduction to Databases
\_ Part - A
| (Short Questions)

\medskip\noindent\textbf{Question 3.} Note : Answer the following (each question carry
| 2 marks) : 10x2=20
: 1. Define the following terms :.
| (a) Data model
| {b} Database schema .
| (c) Client/server architecture
) . (d} Owner entity type
| (ce) Weak entity type
| (f} Entity
REFNO 052213
(gz) Attribute Value "
. (h) Relation state
\begin{parts}
  \item Degree of a relation
  \item Normalization
  PART --- B
  Note: All compulsory : Either or Type Questions.
  5x4=20
\end{parts}

\medskip\noindent\textbf{Question 2.} What is the difference between logical data
independence and physical data
independence ? Which one is harder to
achieve ? Why ?
\centerline{\textbf{OR}}
Discuss the different data types of SQL.

\medskip\noindent\textbf{Question 3.} Discuss the entity integrity and referential
integrity constraints. Why is each considered
important ? .
Discuss Functional dependency with suitable
example.

\medskip\noindent\textbf{Question 4.} What is a relationship type ? Explain the
: differences among a relationship instance, a
relationship type, and a relationship set.
REFNO 052213
\centerline{\textbf{OR}}
What is a participation role ? When is it
necessary to use role names in the description
of relationship types ? : .

\medskip\noindent\textbf{Question 5.} What is meant by a recursive relationship
type ? Give some examples of recursive
relationship types.
\centerline{\textbf{OR}}
Define SQL ? Write the Syntax of create table
and Insert table in the data ? .

\medskip\noindent\textbf{Question 6.} Why are duplicate tuples not allowed in a
relation ?
| OR
| What is the difference between a key and a
| super key ?
: Part - C
i Note: Open choice - 3 out of 4 questions.
| ’ 3x10=30

\medskip\noindent\textbf{Question 7.} Describe the three-schema architecture. Why .
} do we need mappings between scheme
a levels ? How do different schema definition
languages support this architecture ?
Fi
. 3 PLEO:
REFNO 052213
\& | . ’

\medskip\noindent\textbf{Question 8.} Discuss the different type of user-friendly
' interfaces and types of users who typically use
each.

\medskip\noindent\textbf{Question 9.} Discuss the various reasons that lead to the
occurrence of NULL values in relations. |

\medskip\noindent\textbf{Question 10.} Discuss the role of a high-level data model, in
the database design process. ’
KKK ,
4 " 300

\vfill
\rule{\linewidth}{0.4pt}
\begin{center}\small
\href{https://bhu-exam.vercel.app/}{Click here to visit our website}\\[4pt]\href{https://me-aryan.vercel.app/}{Click here to visit our contributor}\\[4pt]\href{https://quashan-me.vercel.app/}{Click here to visit our contributor}
\end{center}
\end{document}
